\documentclass[UTF8,a4paper,11pt]{ctexart}
\usepackage[a4paper,margin=2.1cm,headheight=14pt]{geometry}
\usepackage{booktabs,longtable,tabularx,array}
\usepackage{xcolor}
\usepackage{enumitem}
\usepackage{hyperref}
\usepackage{fancyhdr}
\usepackage{float}
\usepackage{tikz}
\usetikzlibrary{arrows.meta,positioning}
\hypersetup{colorlinks=true,linkcolor=blue,urlcolor=blue}
\definecolor{navy}{HTML}{17313C}
\definecolor{gray}{HTML}{F5F7F8}
\definecolor{bluegray}{HTML}{EAF2FF}
\setlist{nosep,leftmargin=2em}
\setlength{\parindent}{2em}
\setlength{\parskip}{0.35em}
\sloppy
\pagestyle{fancy}
\fancyhf{}
\lhead{企业 OA 业务现状调研}
\rhead{访谈问卷与记录模板}
\cfoot{\thepage}
\newcolumntype{Y}{>{\raggedright\arraybackslash}X}
\newcommand{\question}[1]{\textbf{#1}}
\newcommand{\linefill}{\rule{0.94\linewidth}{0.4pt}}

\title{企业 OA 业务现状调研手册\par\large 以真实订单、单据、数据和岗位协作为核心的访谈问卷}
\author{企业需求访谈、流程还原与范围确认用}
\date{\today}

\begin{document}
\maketitle
\tableofcontents
\newpage

\section{调研目标与使用方法}
\subsection{调研最终要确认的内容}
调研结束后，企业和项目组应共同确认：
\begin{enumerate}
  \item 企业日常最重要的业务对象是什么：客户、合同、订单、项目、采购、生产任务、服务工单或其他对象。
  \item 一笔真实业务从开始到结束经过哪些业务环节，每个环节的完成标准是什么。
  \item 每个环节由哪些岗位参与，谁负责、谁复核、谁批准、谁接收结果。
  \item 每个环节产生或使用哪些单据、票据、表格、附件、邮件和其他资料。
  \item 单据中的字段分别由谁填写，数据从哪里来，哪些是手工填写，哪些是计算或汇总结果。
  \item 数据如何从一个岗位传给下一个岗位，哪些内容被重复抄录，哪些内容容易遗漏或出错。
  \item 管理人员、业务人员、财务人员分别需要看哪些信息，并据此做什么决定。
  \item 业务发生变更、退回、延期、取消、超预算、部分交付或争议时，企业实际如何处理。
  \item 哪些流程和资料必须纳入系统，哪些继续使用现有方式即可，哪些暂时不应纳入。
\end{enumerate}

\section{第一部分：企业业务画像}
\subsection{企业与业务模式}
\begin{longtable}{p{0.31\textwidth}p{0.59\textwidth}}
\toprule
问题 & 记录 \\
\midrule
企业主要做什么产品或服务？ & \\
主要客户是谁？客户如何提出需求？ & \\
业务是一次性订单、长期合同、项目制、生产制、服务制还是混合？ & \\
一项业务通常持续多久？是否存在跨月、跨季度、跨年度业务？ & \\
每月大约有多少客户、订单、项目、采购和结算事项？ & \\
哪些业务类型的处理方式不同？例如大客户、小客户、定制业务、标准业务。 & \\
业务中最关键的三个时间点是什么？ & \\
目前最容易延误、出错、漏办或无法追责的环节是什么？ & \\
管理层最希望及时知道哪些情况？ & \\
\bottomrule
\end{longtable}

\subsection{企业对“OA”的实际期待}
不要接受“想做一个完整 OA”作为最终答案，继续询问：
\begin{enumerate}
  \item 企业说的 OA 主要是审批和协作，还是要管理订单/项目的全过程？
  \item 当前最急需改善的是哪一项具体业务？为什么现在必须解决？
  \item 如果只能先解决一个流程，选择哪个流程？完成后谁会直接受益？
  \item 哪些工作目前已经由其他系统稳定处理，不希望重复处理？
  \item 哪些工作虽然存在问题，但由于频率低、规则不清或责任人不明确，暂时不适合系统化？
  \item 企业希望系统帮助“记录发生过什么”，还是帮助“控制下一步能不能做”，还是帮助“汇总后进行管理决策”？
\end{enumerate}

\section{第二部分：用真实业务还原全过程}
\subsection{真实业务样本}
\begin{longtable}{p{0.31\textwidth}p{0.59\textwidth}}
\toprule
项目 & 记录 \\
\midrule
业务编号、客户/合作方、业务名称 & \\
业务类型和来源 & \\
开始日期、计划结束日期、实际结束日期 & \\
金额、数量、规模、主要交付内容 & \\
当前状态或最终结果 & \\
涉及部门和岗位 & \\
涉及的全部单据、票据和附件 & \\
本案例中发生过的延期、退回、变更、争议或补录 & \\
\bottomrule
\end{longtable}

\subsection{全过程记录表}
请按业务实际发生顺序填写。不要用软件菜单名称代替业务环节名称。
\begin{longtable}{p{0.07\textwidth}p{0.16\textwidth}p{0.17\textwidth}p{0.17\textwidth}p{0.19\textwidth}p{0.16\textwidth}}
\toprule
序号 & 业务环节 & 谁发起/谁负责 & 进入条件 & 处理内容与完成标准 & 交给谁/产生什么资料 \\
\midrule
1 & & & & & \\
2 & & & & & \\
3 & & & & & \\
4 & & & & & \\
5 & & & & & \\
6 & & & & & \\
7 & & & & & \\
8 & & & & & \\
9 & & & & & \\
10 & & & & & \\
\bottomrule
\end{longtable}

\subsection{每个业务环节都要问}
\begin{enumerate}
  \item 这个环节在企业内部叫什么？不同部门是否有不同叫法？
  \item 是什么事情触发了这个环节？谁提出、谁通知、谁决定开始？
  \item 处理人开始工作时，手上必须有哪几份资料？资料从谁那里来？
  \item 处理人具体做了哪些工作？哪些是判断，哪些是填写，哪些是计算，哪些是沟通？
  \item 什么结果才算完成？由谁确认完成？
  \item 完成后产生哪些表单、票据、附件、编号、通知或待办事项？
  \item 下一岗位拿到什么内容？是否需要重新录入或重新核对？
  \item 哪些内容一旦确认就不能修改？如果必须修改，谁同意？
  \item 该环节通常耗时多久？等待时间发生在哪里？
  \item 该环节最常见的错误、遗漏和争议是什么？
  \item 这个环节的数据后续被谁查看、用于什么工作或决定？
\end{enumerate}

\section{第三部分：订单、项目和业务状态}
\subsection{企业如何称呼和区分业务对象}
\begin{longtable}{p{0.22\textwidth}p{0.19\textwidth}p{0.24\textwidth}p{0.27\textwidth}}
\toprule
对象 & 企业名称/编号 & 与其他对象的关系 & 从开始到结束的状态 \\
\midrule
客户/合作方 & & & \\
客户需求/商机 & & & \\
报价/方案 & & & \\
合同 & & & \\
订单/项目 & & & \\
采购申请/采购单 & & & \\
交付任务/服务事项 & & & \\
验收/结算 & & & \\
发票/收款 & & & \\
费用/报销 & & & \\
\bottomrule
\end{longtable}

\subsection{状态确认问题}
对订单、项目、合同、采购、费用和验收分别询问：
\begin{itemize}
  \item 企业实际使用哪些状态名称？每个状态是什么意思？
  \item 进入该状态的条件是什么？谁可以让它进入该状态？
  \item 该状态下允许做什么，不允许做什么？
  \item 什么情况下退回、暂停、取消、作废、重新打开或重新开始？
  \item 状态改变后，谁会收到通知？谁必须继续处理？
  \item 状态是否代表业务事实，还是只代表某个人已经处理过？
  \item 一个业务是否可能同时处于多个状态，例如“已交付但未结算”“已开票但未收款”？
  \item 企业是否存在“看似完成但实际上还缺资料”的状态？目前如何识别？
\end{itemize}

\section{第四部分：单据、票据、表格和附件目录}
\subsection{全部资料盘点表}
“单据”包括纸质文件、Excel、Word、PDF、邮件、聊天截图、照片、扫描件和系统内记录。
\begin{longtable}{p{0.17\textwidth}p{0.17\textwidth}p{0.18\textwidth}p{0.21\textwidth}p{0.17\textwidth}}
\toprule
资料名称/俗称 & 类型 & 在哪个环节产生 & 谁填写/确认/使用 & 保存在哪里 \\
\midrule
& 报价/合同/申请/审批/采购/交付/财务/其他 & & & \\
& & & & \\
& & & & \\
& & & & \\
& & & & \\
& & & & \\
& & & & \\
& & & & \\
\bottomrule
\end{longtable}

\subsection{一张单据的详细调查表}
\begin{longtable}{p{0.32\textwidth}p{0.58\textwidth}}
\toprule
调查问题 & 企业实际情况 \\
\midrule
单据正式名称、内部俗称、编号规则 & \\
它要证明什么或推动什么业务？ & \\
由哪个环节产生，前置资料是什么？ & \\
谁填写、谁复核、谁审批、谁接收？ & \\
纸质、Excel、Word、PDF、邮件、聊天还是其他方式？ & \\
是否有固定模板？模板由谁维护？ & \\
每个字段的名称、含义、填写示例 & \\
哪些字段必须填，哪些字段可以空，哪些字段在特定情况下才必填？ & \\
字段是手工填写、从其他资料抄录，还是由企业规则计算？ & \\
金额、数量、日期、税额、折扣等字段的口径是什么？ & \\
需要哪些附件、原件、签字、盖章或客户确认？ & \\
谁根据这张单据决定下一步？依据是什么？ & \\
单据通过后是否允许修改？修改要不要重新确认？ & \\
退回、作废、重开、补录时使用什么方式？ & \\
最终由谁查询、对账、汇总或归档？ & \\
需要保存多久，原件在哪里，查找是否困难？ & \\
\bottomrule
\end{longtable}

\subsection{票据专项调查}
对合同、报价单、发票、付款凭证、收款凭证、验收单、结算单等逐项确认：
\begin{enumerate}
  \item 它是内部管理资料，还是对客户、供应商、财务、税务或审计有效的正式凭证？
  \item 原件由谁保管？企业是否需要保存扫描件、照片或复印件？
  \item 是否需要签字、盖章、电子签名或客户确认邮件？
  \item 一张票据对应一笔业务，还是可能对应多个订单/项目？
  \item 是否存在分批开具、部分支付、冲销、作废、补开或重复提交？
  \item 金额是否区分含税、未税、税额、折扣、预付款、尾款和扣款？
  \item 谁核对票据与订单、合同、收货、验收和付款是否一致？
  \item 票据缺失时业务能否继续？谁批准例外？
\end{enumerate}

\section{第五部分：字段与数据来源调查}
\subsection{业务字段清单}
对每张核心单据和每个核心业务对象填写，不要只列字段名称。
\begin{longtable}{p{0.16\textwidth}p{0.20\textwidth}p{0.15\textwidth}p{0.18\textwidth}p{0.21\textwidth}}
\toprule
字段名称 & 企业含义/填写示例 & 谁提供 & 什么时候填写/修改 & 谁查看/使用及用途 \\
\midrule
编号 & & & & \\
客户/合作方 & & & & \\
负责人/部门 & & & & \\
业务内容 & & & & \\
数量/单位 & & & & \\
金额/币种/税额 & & & & \\
日期/期限 & & & & \\
状态/结果 & & & & \\
审批/确认信息 & & & & \\
附件/凭证 & & & & \\
备注/特殊要求 & & & & \\
\bottomrule
\end{longtable}

\subsection{逐字段追问}
\begin{enumerate}
  \item 这个字段在企业内部是什么意思？有没有同名不同义的情况？
  \item 谁最先知道或提供这个信息？谁负责把它写下来？
  \item 信息来自客户、合同、报价、现场、供应商、财务还是其他人员？
  \item 是一次填写后不变，还是会随着业务推进而变化？
  \item 谁可以修改？修改后是否需要告诉其他人？
  \item 如果没有这个信息，业务能否继续？
  \item 该字段是否有固定格式、编号规则、单位、上下限或可选值？
  \item 是否需要与其他单据中的同一信息核对？以哪一份为准？
  \item 该字段是否用于计算、审批、提醒、对账、排名或管理报告？
  \item 历史资料中是否有这个字段？如果没有，企业如何补齐？
\end{enumerate}

\section{第六部分：岗位分工与数据使用}
\subsection{岗位协作表}
\begin{longtable}{p{0.17\textwidth}p{0.19\textwidth}p{0.18\textwidth}p{0.20\textwidth}p{0.18\textwidth}}
\toprule
岗位/部门 & 负责哪些环节 & 要填写哪些资料 & 要查看哪些资料 & 根据什么结果做决定 \\
\midrule
业务/销售 & & & & \\
项目负责人 & & & & \\
采购 & & & & \\
仓储/交付 & & & & \\
财务 & & & & \\
部门负责人 & & & & \\
管理层 & & & & \\
客户/供应商 & & & & \\
\bottomrule
\end{longtable}

\subsection{谁需要什么信息}
对每类使用者询问：
\begin{itemize}
  \item 他在什么时间查看信息？每天、每周、月底、审批前还是出现问题时？
  \item 他需要看到单笔业务详情，还是只需要汇总结果？
  \item 他最关心哪些字段：进度、金额、成本、利润、欠款、风险、责任人、期限还是凭证？
  \item 他看完数据后要做什么动作：批准、催办、安排资源、付款、开票、追款、调整计划还是追责？
  \item 他是否需要看到原始单据和附件？是否需要导出、打印或转发？
  \item 哪些信息不能让其他岗位看到？哪些信息必须让所有相关人员看到？
\end{itemize}

\subsection{责任确认}
\begin{longtable}{p{0.25\textwidth}p{0.20\textwidth}p{0.20\textwidth}p{0.25\textwidth}}
\toprule
事项 & 实际负责人 & 最终批准人 & 负责人不在时如何处理 \\
\midrule
创建业务记录 & & & \\
确认客户需求/报价 & & & \\
确认合同/订单 & & & \\
发起采购 & & & \\
审核费用/票据 & & & \\
确认交付/验收 & & & \\
项目结算 & & & \\
开票/收款跟进 & & & \\
异常和争议处理 & & & \\
\bottomrule
\end{longtable}

\section{第七部分：数据传递、重复录入与汇总}
\subsection{数据流转调查}
\begin{longtable}{p{0.19\textwidth}p{0.19\textwidth}p{0.22\textwidth}p{0.20\textwidth}p{0.14\textwidth}}
\toprule
数据/资料 & 从谁到谁 & 当前通过什么方式传递 & 接收方如何核对 & 是否重复填写 \\
\midrule
客户需求 & & & & \\
合同/订单信息 & & & & \\
报价与预算 & & & & \\
采购和供应商信息 & & & & \\
进度和交付结果 & & & & \\
费用和票据 & & & & \\
验收和结算 & & & & \\
开票和收款 & & & & \\
\bottomrule
\end{longtable}

\subsection{汇总表和管理报表}
\begin{longtable}{p{0.22\textwidth}p{0.18\textwidth}p{0.21\textwidth}p{0.21\textwidth}p{0.12\textwidth}}
\toprule
报表/汇总名称 & 谁制作/谁看 & 包含哪些内容 & 数据来自哪些单据 & 使用频率 \\
\midrule
订单/项目进度表 & & & & \\
销售/业务统计 & & & & \\
采购和供应商汇总 & & & & \\
费用/报销汇总 & & & & \\
合同/开票/收款汇总 & & & & \\
成本/利润汇总 & & & & \\
逾期/异常事项表 & & & & \\
管理层经营汇报 & & & & \\
\bottomrule
\end{longtable}

追问重点：目前谁花时间汇总？每次需要多久？是否从多个 Excel 复制？是否反复向业务人员要数据？不同报表的数字为什么可能不一致？管理层看到报表后会采取什么行动？

\section{第八部分：异常和非标准业务}
正常流程之外的情况往往决定企业真正的管理需求。请结合真实案例填写：
\begin{longtable}{p{0.25\textwidth}p{0.28\textwidth}p{0.35\textwidth}}
\toprule
异常场景 & 企业当前怎么处理 & 需要保留的资料、责任和结果 \\
\midrule
客户临时改需求 & & \\
合同金额或范围变更 & & \\
报价反复修改或已确认后重谈 & & \\
订单拆分、合并或转给其他人 & & \\
业务延期、暂停或取消 & & \\
采购超预算或紧急采购 & & \\
供应商更换或分批交付 & & \\
没有发票或票据不合格 & & \\
员工垫付和重复报销 & & \\
部分交付或部分验收 & & \\
客户拒收、投诉或验收争议 & & \\
部分开票、部分回款或逾期 & & \\
跨部门/跨公司分摊 & & \\
资料填错或审批后发现错误 & & \\
负责人离职或岗位交接 & & \\
\bottomrule
\end{longtable}

每个异常都要问：谁发现、谁决定、谁批准、用什么凭证证明、原来的数据是否保留、下一步如何继续、最后如何关闭。

\section{第九部分：现有工作方式与已有工具}
本部分只调查企业业务事实，不评价软件好坏。
\begin{longtable}{p{0.20\textwidth}p{0.22\textwidth}p{0.25\textwidth}p{0.22\textwidth}}
\toprule
工作内容 & 现在用什么方式处理 & 哪些岗位参与 & 当前最不满意之处 \\
\midrule
客户/订单登记 & & & \\
合同和报价 & & & \\
项目/任务推进 & & & \\
采购和供应商沟通 & & & \\
费用和票据 & & & \\
交付和验收 & & & \\
开票和收款 & & & \\
报表和经营汇报 & & & \\
通知、催办和交接 & & & \\
资料归档和查询 & & & \\
\bottomrule
\end{longtable}

\subsection{现有方式的深入追问}
\begin{itemize}
  \item 现在的做法为什么形成？哪些是企业规定，哪些只是个人习惯？
  \item 哪些表格是必须保留的正式资料，哪些只是个人辅助记录？
  \item 如果负责人员请假或离职，其他人能否接手？资料是否找得到？
  \item 哪些信息只能在某个人的电脑、微信或聊天记录里找到？
  \item 企业现在是否已经有一套大家认可的客户、订单、金额、项目状态和供应商名单？
  \item 当前方式中最费时间的步骤是什么？最容易错的步骤是什么？最难追责的步骤是什么？
\end{itemize}

\section{第十部分：范围判断，不预设必须做完整 OA}
\subsection{业务模块评估表}
请根据前面收集到的事实，不根据想象评分。每项 1--5 分。
\begin{longtable}{p{0.22\textwidth}p{0.13\textwidth}p{0.13\textwidth}p{0.13\textwidth}p{0.29\textwidth}}
\toprule
业务领域 & 发生频率 & 当前痛点 & 规则清晰度 & 是否纳入一期及理由 \\
\midrule
客户与业务机会 & & & & \\
订单/项目推进 & & & & \\
报价与合同 & & & & \\
采购与供应商 & & & & \\
交付与验收 & & & & \\
费用与报销 & & & & \\
开票与收款 & & & & \\
库存/物料/生产 & & & & \\
行政审批 & & & & \\
人事相关事务 & & & & \\
经营报表 & & & & \\
\bottomrule
\end{longtable}

\subsection{判断题}
\begin{longtable}{p{0.48\textwidth}p{0.40\textwidth}}
\toprule
问题 & 结论和依据 \\
\midrule
是否有一条跨部门、重复发生、责任清晰的核心业务链？ & \\
该业务链是否存在大量重复登记、人工催办或信息丢失？ & \\
企业是否能明确这条业务链的开始、结束和完成标准？ & \\
涉及岗位是否愿意共同使用同一套资料和状态？ & \\
已有财务/ERP 等系统是否已覆盖其中一部分？ & \\
企业需要的是过程协作、业务记录、管理汇总，还是三者都有？ & \\
是否有一个部门负责人对一期结果负责？ & \\
一期是否能提供真实样例、历史资料和人员配合？ & \\
最适合一期的业务闭环是什么？ & \\
暂不做哪些内容？为什么？ & \\
\bottomrule
\end{longtable}

\subsection{三种常见结论}
\begin{description}
  \item[完整 OA] 企业有多类稳定且相互关联的跨部门流程，并且需要统一的日常协作、审批、资料和管理入口。
  \item[业务闭环系统] 企业真正痛点集中在订单/项目、采购、费用、交付、合同或回款中的一条主线，先解决这条主线更容易取得结果。
  \item[现状梳理项目] 企业各部门对流程、字段和责任的理解尚不一致，当前最需要的是统一业务规则和资料，而不是立即建设完整系统。
\end{description}

\section{第十一部分：调研结论和企业确认}
\subsection{企业业务结论}
\begin{longtable}{p{0.30\textwidth}p{0.60\textwidth}}
\toprule
结论项 & 记录 \\
\midrule
企业最核心的业务对象 & \\
最重要的业务流程起点和终点 & \\
一期要解决的最主要问题 & \\
涉及的业务部门和岗位 & \\
必须保留的单据、票据和附件 & \\
必须确认的关键字段和金额口径 & \\
必须看到数据的人员及使用目的 & \\
当前最严重的异常或管理风险 & \\
已经明确的业务规则 & \\
不同部门之间尚未统一的说法 & \\
待企业补充的资料 & \\
待企业负责人决定的事项 & \\
\bottomrule
\end{longtable}

\subsection{纳入与不纳入清单}
\begin{longtable}{p{0.17\textwidth}p{0.25\textwidth}p{0.45\textwidth}}
\toprule
范围 & 业务内容 & 原因/依据/后续安排 \\
\midrule
一期纳入 & & \\
一期纳入 & & \\
后续考虑 & & \\
后续考虑 & & \\
明确不做 & & \\
明确不做 & & \\
\bottomrule
\end{longtable}

\subsection{必须形成的调研资料}
\begin{enumerate}
  \item 真实业务全过程说明；
  \item 业务环节和岗位责任表；
  \item 单据、票据、表格和附件目录；
  \item 核心单据字段清单及字段来源；
  \item 数据传递和人工汇总说明；
  \item 角色与信息使用清单；
  \item 异常处理案例；
  \item 企业术语和状态名称表；
  \item 一期业务范围及暂缓范围；
  \item 待确认问题、负责人和截止时间。
\end{enumerate}

\section{访谈结束时的最小必问清单}
时间有限时，以下问题不能省略：
\begin{enumerate}
  \item 请拿出一笔真实订单/项目，从开始到结束讲一遍。
  \item 每一步由谁负责，开始时需要什么资料，完成后交给谁？
  \item 这笔业务产生了哪些表单、票据、附件、Excel、邮件和聊天记录？
  \item 每张核心单据有哪些字段？每个字段是谁填、从哪里来、谁使用？
  \item 哪些资料或数字会被重复抄到其他表里？谁负责汇总？
  \item 业务的金额、数量、成本、开票、收款和结算分别以什么为准？
  \item 谁需要看单笔业务，谁需要看汇总结果，分别用来做什么？
  \item 业务退回、变更、延期、取消、超预算、缺票或发生争议时怎么处理？
  \item 当前最耗时、最容易错、最难找到资料、最难追责的三个地方是什么？
  \item 如果一期只改善一个业务环节，企业认为哪个结果最有价值？
\end{enumerate}

\end{document}
